%%% SECTION 5: ADAPTED 21 CFR PART 50 FOR PHYSICAL AI PROTECTION %%%

The Physical AI Adaption: 21 CFR Part 50 \cite{cfr50-adapt} provides
comprehensive human subjects protection specifically designed for Physical AI
oncology trials, building on the existing federal regulation \cite{cfr-part-50}.
By adapting a well-established regulation rather than proposing new legislation,
this work adds immediate credibility and practical applicability. A critical
theme throughout this adaptation is that control has shifted towards the
patient's side, with expanded patient rights including new protections specific
to robotic and AI-assisted medical procedures that have no precedent in
traditional clinical trial regulation.

\subsection{Scope and Physical AI Definitions}
\label{subsec:cfr50-scope}

The adapted Part 50 applies to all clinical investigations regulated by FDA that
involve Physical AI systems interacting with human subjects. Subpart A (General
Provisions) establishes the scope and provides 17 Physical AI-specific
definitions that standardize terminology for robotic and AI-assisted clinical
procedures.

Key definitions include: \textit{Physical AI System}, encompassing any robotic
platform integrated with artificial intelligence that interacts directly with
trial participants; \textit{Physical AI Adverse Event}, any untoward occurrence
during a clinical investigation that is attributable to Physical AI system
operation; \textit{Physical AI Operator}, a qualified individual who monitors
and controls Physical AI systems during clinical procedures; \textit{Robotic
Procedure}, any clinical intervention performed wholly or partly by a Physical
AI system; \textit{Autonomous Mode}, operation of a Physical AI system where AI
algorithms direct robotic actions without real-time human input;
\textit{Supervised Mode}, operation where human operators maintain continuous
control; and \textit{Physical AI Safety Matrix}, a pre-procedure verification
checklist ensuring all safety systems are operational.

These definitions create the foundation for all subsequent protections in the
regulation. The distinction between autonomous and supervised modes is
particularly important because it determines the level of informed consent
detail required and the intensity of ongoing monitoring. Cross-referencing with
the 22 definitions in the adapted 21 CFR Part 312 \cite{cfr312-adapt} and
the 30 definitions in the adapted ICH E6(R3) \cite{ich-adapt} ensures
terminological consistency across the entire regulatory framework.

\subsection{Informed Consent for Physical AI Procedures}
\label{subsec:cfr50-consent}

Subpart B adapts the general requirements for informed consent to Physical AI
contexts. Traditional informed consent requires disclosure of trial purpose,
procedures, risks, benefits, alternatives, and confidentiality protections. The
Physical AI adaptation adds specific consent elements required for any trial
involving robotic systems.

The additional consent elements include: a description of each Physical AI
system type and its role in the participant's care, including the robot
manufacturer, model, and general capabilities; an explanation of how AI
algorithms participate in treatment planning, monitoring, and decision-making;
a clear statement of the participant's right to request human-only care at any
point during the trial, including during active robotic procedures; a
description of data collection specific to Physical AI systems, including robot
telemetry, video recording, sensor data, and AI decision logs, with associated
privacy protections; an explanation of emergency stop procedures and how
participants can signal distress or discomfort during robotic interactions; and
a description of the safety monitoring systems in place, including who monitors
robot performance and how deviations are detected and addressed.

These consent elements represent a fundamental shift in the informed consent
paradigm. In traditional trials, consent primarily addresses drug or device
risks and benefits. In Physical AI trials, consent must address a new category
of risk: the interaction between a human patient and an autonomous or
semi-autonomous robotic system that physically touches, manipulates, or
occupies space near the patient. Patient Instructions: Physical AI Trials
\cite{patient-instructions-paper} provides the patient-facing implementation
of these consent requirements, translating regulatory language into accessible
information for each of the ten robot types.

\subsection{Physical AI System Safety Requirements}
\label{subsec:cfr50-safety}

Subpart C is a new subpart created specifically for Physical AI trials,
establishing additional protections for subjects in Physical AI clinical
investigations. This subpart has no precedent in the original 21 CFR Part 50
and represents the most significant structural addition to the adapted
regulation.

\subsubsection{Pre-Procedure Safety Matrix}

Before any Physical AI procedure begins, the Physical AI Safety Matrix must be
completed and documented. This matrix verifies: robot system power and
mechanical status (all joints responsive, actuators calibrated, end-effector
functional); sensor calibration status (force sensors, position sensors,
proximity sensors, temperature sensors all within specification); emergency stop
system verification (hardware E-stop functional, software E-stop responsive,
wireless E-stop paired); communication system verification (robot-to-server
connection active, data logging operational, alert system functional); patient
identification confirmation (biometric or badge-based patient ID matched to
protocol records); workspace clearance verification (operating zone free of
unauthorized objects, personnel positioned safely, barriers in place); and AI
model status verification (approved model version loaded, performance within
specification, no pending updates requiring validation).

\subsubsection{Runtime Safety Monitoring}

During Physical AI procedures, continuous monitoring requirements include:
real-time force and torque monitoring at all patient contact points, with
automatic procedure pause when force exceeds protocol-specified limits;
position deviation monitoring comparing actual robot position to planned
trajectory, with alerts when deviation exceeds tolerance; vital sign
integration connecting patient monitoring systems to Physical AI control
systems, enabling automatic response to physiological changes; AI decision
monitoring logging all autonomous decisions with confidence scores and
alternative actions considered; and video recording of all patient-robot
interactions for post-procedure review and adverse event investigation.

\subsubsection{Post-Procedure Requirements}

After Physical AI procedures, required documentation includes: a complete
Physical AI procedure report recording all robot actions, parameters, and
outcomes; patient follow-up specific to Physical AI interactions, including
assessment of any physical marks, discomfort, or psychological responses to
robotic interaction; data archival ensuring all telemetry, sensor, video, and
decision log data are properly stored with appropriate retention periods; and
system reset verification confirming the Physical AI system is properly reset
and ready for the next procedure.

\subsubsection{Task-Order Lifecycle}

The adaptation establishes a formal lifecycle for Physical AI task orders from
initiation through completion. Each task order must be authorized by a
qualified Physical AI operator, with the authorization level matched to the
task complexity and risk. Task execution must be monitored against approved
parameters, with any deviation triggering documentation and potential
escalation. Task completion must be verified through both automated system
checks and human confirmation before the system is released for the next task.

\subsubsection{Forbidden Operations}

The adapted regulation explicitly prohibits Physical AI systems from performing
certain operations without human authorization: initiating surgical incisions
or tissue removal without real-time surgeon oversight; modifying medication
dosages or administration schedules without pharmacist or physician
confirmation; altering treatment plans based solely on AI recommendations
without clinical team review; performing emergency medical interventions
without qualified human medical staff present; and accessing or modifying
patient records beyond the specific data required for the current authorized
task.

\subsection{IRB Review of Physical AI Investigations}
\label{subsec:cfr50-irb}

The adapted regulation establishes specific requirements for IRB review of
Physical AI trials that extend beyond traditional IRB evaluation criteria.
IRBs must evaluate: the technical qualifications of Physical AI system
operators and maintenance staff; the adequacy of Physical AI system validation
evidence, including simulation results and bench testing data; the
appropriateness of Physical AI safety monitoring plans; the completeness of
Physical AI-specific informed consent elements; and the adequacy of contingency
plans for Physical AI system failures.

Ongoing monitoring requirements specify that IRBs must receive periodic reports
on Physical AI system performance, including uptime, adverse events, near-miss
incidents, and maintenance activities. Protocol modification triggers specific
to Physical AI include: Physical AI system software updates that change
operational parameters, introduction of new Physical AI system types not covered
in the original protocol, changes to Physical AI safety monitoring procedures,
and changes to the degree of autonomy granted to Physical AI systems. The NCI
CIRB \cite{nci-cirb} provides a model for centralized Physical AI review that
could enable consistent evaluation across multiple trial sites.

\subsection{Ongoing Consent and Subject Notification}
\label{subsec:cfr50-ongoing}

Unlike traditional trials where informed consent is primarily a one-time event,
Physical AI trials require ongoing consent processes because the technology may
be updated, new robots may be introduced, and AI algorithms may be refined
during the trial period. Participants must be informed about: Physical AI system
software updates that change system behavior or capabilities; introduction of
new robot types at the trial site; changes to safety monitoring procedures;
results of periodic Physical AI system safety reviews; and any Physical AI
adverse events that occurred at the trial site that may affect the
participant's risk assessment.

This ongoing consent process gives patients continuous agency over their
participation. A patient who consented to interact with a da Vinci surgical
robot may reconsider if a humanoid robot is later added to the trial site for
logistics support. A patient who accepted a specific AI algorithm for treatment
planning may want to reconsider if the algorithm is updated. The adapted
regulation ensures that patients are never surprised by changes to the Physical
AI systems involved in their care.

\subsection{Data Protection and Privacy}
\label{subsec:cfr50-data}

The adapted regulation establishes protections for the additional categories of
data generated by Physical AI systems. Robot telemetry data recording position,
force, speed, and environmental readings constitutes a new category of health
information that was not contemplated by the original Part 50. Video recordings
of Physical AI procedures create visual records of patient interactions with
robotic systems. Sensor data from patient-robot interactions may include
biometric measurements beyond those collected in traditional trials. AI decision
logs that document the reasoning behind automated treatment decisions create
records that may have implications for patient rights and legal accountability.

All these data categories are subject to HIPAA \cite{hipaa} protections when
they constitute protected health information. The federated learning framework
\cite{fl-paper} and the broader literature on federated learning in healthcare
\cite{federated-learning-healthcare} provide approaches for utilizing these
data for model improvement while maintaining patient privacy through
decentralized processing.

\subsection{Physical AI System Classification}
\label{subsec:cfr50-classification}

The adapted regulation establishes a classification framework organizing
Physical AI systems by risk level into three tiers. Tier 1 (Standard Risk)
encompasses Physical AI systems with minimal direct patient contact, such as
logistics robots and specimen handling cobots, which require standard informed
consent and monitoring. Tier 2 (Elevated Risk) encompasses systems with
moderate patient interaction, such as positioning robots and imaging robots,
requiring enhanced consent elements and continuous monitoring. Tier 3 (High
Risk) encompasses systems with direct therapeutic interaction, such as surgical
robots and needle-placement robots, requiring comprehensive consent with
procedure-specific elements, real-time safety monitoring, and qualified
operator presence. The classification tier determines the regulatory pathway,
monitoring intensity, and consent detail required, connecting to the USL
scoring framework in Section~\ref{sec:psl-usl} for quantitative assessment.

\subsection{Pediatric Protections in Physical AI Trials}
\label{subsec:cfr50-pediatric}

Subpart D adapts existing pediatric protections for Physical AI contexts.
Modified consent and assent procedures address minors interacting with robots,
requiring age-appropriate explanations of robotic systems. Additional safety
requirements for pediatric Physical AI procedures include reduced force limits,
smaller workspace boundaries, and enhanced human oversight. IRB duties specific
to pediatric Physical AI investigations include evaluation of whether robotic
interactions may cause additional psychological distress in pediatric patients.
Pediatric oncology is a particularly important application where Physical AI may
reduce treatment burden on young patients through faster procedures, reduced
need for repeated positioning, and companion robot support during treatment.

\subsection{Glossary of Physical AI Terms}
\label{subsec:cfr50-glossary}

The adapted Part 50 concludes with a supplementary glossary of 35 Physical AI
terms that complement the regulatory definitions in Subpart A. Terms include
Agentic AI (AI systems capable of autonomous multi-step task execution),
Calibration (the process of adjusting Physical AI system parameters to meet
specified accuracy standards), Digital Twin (a computational model replicating a
physical patient or system for simulation and prediction), MCP Server (a Model
Context Protocol server providing standardized AI-tool communication), and USL
(Unification Standard Level, a quantitative robot readiness scoring framework).
These terms are cross-referenced with the ICH glossary
(Section~\ref{subsec:ich-glossary}) to ensure consistency across all adapted
standards in the National Platform.

\subsection{Three-Tier Physical AI Classification System}
\label{subsec:cfr50-three-tier}

The three-tier classification introduced in
Section~\ref{subsec:cfr50-classification} warrants detailed examination because
it determines the regulatory pathway, monitoring intensity, consent complexity,
and reporting obligations for every Physical AI system deployed in an oncology
clinical trial. Table~\ref{tab:three-tier-detail} presents the complete
classification framework with representative robot types, regulatory
requirements, and consent elements for each tier.

\begin{table}[H]
\centering
\caption{Three-Tier Physical AI Classification System: Requirements and Consent Elements}
\label{tab:three-tier-detail}
\small
\begin{tabularx}{\textwidth}{l X X X}
\toprule
\textbf{Attribute} & \textbf{Tier 1 - Standard Risk} & \textbf{Tier 2 - Elevated Risk} & \textbf{Tier 3 - High Risk} \\
\midrule
Representative robots & Logistics humanoids (Digit, Optimus), social companion robots, specimen-handling cobots & RT positioning robots, imaging robots, rehabilitation exoskeletons, RT motion-tracking robots & Surgical robots (da Vinci, Hugo RAS, Versius), needle-placement robots, steerable needle robots \\
\midrule
Patient contact level & Minimal or no direct physical contact; robot operates in shared space & Moderate direct contact; robot touches or repositions patient under supervision & Direct therapeutic contact; robot performs invasive or semi-invasive procedures \\
\midrule
Autonomy permitted & Supervised or semi-autonomous for logistics tasks; autonomous navigation in designated corridors & Supervised mode required during patient contact; autonomous calibration permitted & Supervised mode mandatory at all times; autonomous incision or tissue manipulation prohibited \\
\midrule
Minimum USL band & Foundational (3.0 or higher) & Intermediate (5.0 or higher) & Advanced (7.0 or higher) \\
\midrule
Operator requirement & Trained staff member within radio or network contact & Certified Physical AI operator present in the room & Board-certified physician plus certified Physical AI operator at console \\
\midrule
Safety matrix scope & Abbreviated matrix covering power, communication, and identification checks & Standard matrix covering all seven verification categories & Extended matrix with procedure-specific additions, dual-operator sign-off, and anesthesia integration \\
\midrule
Monitoring during operation & Periodic status checks via dashboard at 15-minute intervals & Continuous telemetry with automated alerts for parameter exceedance & Real-time multi-channel monitoring including force, position, vital signs, and AI decision confidence \\
\midrule
Adverse event reporting & Category B or C timelines (15 calendar days) for most events & Category A timelines (7 calendar days serious, 15 days non-serious) for contact events & Category A expedited timelines (7 calendar days) for all events; 24-hour reporting for serious surgical events \\
\midrule
IRB review depth & Standard review with Physical AI addendum & Enhanced review with Physical AI safety plan evaluation & Full board review with independent Physical AI technical expert consultation \\
\midrule
Consent elements & General Physical AI disclosure, right to refuse robotic logistics, data collection notice & All Tier 1 elements plus robot-specific interaction description, positioning procedure walkthrough, sensor data explanation & All Tier 1 and 2 elements plus surgical procedure details, autonomous decision-making disclosure, emergency conversion plan, tissue interaction risks \\
\midrule
Post-procedure documentation & System log summary and patient feedback form & Complete telemetry archive, patient assessment, and 24-hour follow-up & Full procedure report, video archive, force-torque analysis, pathology correlation, and 72-hour follow-up \\
\midrule
Re-consent triggers & New robot type added to logistics fleet & Software update changing positioning algorithm; new imaging modality introduced & Any change to surgical robot model, software version, or autonomy level \\
\bottomrule
\end{tabularx}
\end{table}

The tier system creates proportional regulatory burden: Tier 1 systems that
pose minimal risk to patients face streamlined requirements, while Tier 3
systems performing invasive procedures face the most rigorous oversight. This
proportionality is essential for practical deployment because applying Tier 3
requirements to every logistics robot would create unnecessary burden without
improving patient safety, while applying Tier 1 requirements to surgical robots
would provide insufficient protection. The tier classification is determined at
the protocol level during IND submission and reviewed by the IRB as part of the
Physical AI safety plan evaluation.

When a robot's function spans multiple tiers, the higher tier governs. For
example, a cobot that both handles specimens (Tier 1) and assists with patient
positioning (Tier 2) is classified as Tier 2 for all regulatory purposes. This
conservative approach ensures that no Physical AI interaction falls below the
protection level warranted by its highest-risk function.

The tier system also governs training requirements for site personnel. Tier 1
operations require general Physical AI awareness training (minimum 8 hours)
covering robot identification, basic safety procedures, and emergency stop
activation. Tier 2 operations require specialized Physical AI operator
certification (minimum 40 hours) covering robot-specific operation, patient
interaction protocols, and troubleshooting. Tier 3 operations require advanced
Physical AI surgical team certification (minimum 80 hours plus proctored cases)
covering console operation, instrument exchange, emergency conversion, and
team communication protocols.

\subsection{Comprehensive Informed Consent Elements for Physical AI Trials}
\label{subsec:cfr50-consent-comprehensive}

Section~\ref{subsec:cfr50-consent} introduced the Physical AI-specific consent
elements required by the adapted Part 50. The following comprehensive itemized
list specifies every consent element that must be addressed in the informed
consent document for Physical AI oncology trials, organized by category.

\subsubsection{General Physical AI Disclosure Elements}

All Physical AI trial consent forms must include the following general
disclosure elements regardless of tier classification:

\begin{itemize}
    \item A statement that the clinical trial involves Physical AI systems
    (robots with integrated artificial intelligence) that will participate in
    some or all aspects of the participant's care during the trial period.

    \item An enumeration of every Physical AI system type that the participant
    may encounter during the trial, identified by category (surgical robot,
    cobot, humanoid, companion robot, etc.), manufacturer name, model
    designation, and general functional description.

    \item A plain-language explanation of how artificial intelligence
    participates in treatment planning, monitoring, decision support, and
    operational coordination, including a statement about the role of human
    oversight in AI-assisted decisions.

    \item A description of the Physical AI safety monitoring infrastructure,
    including who monitors robot performance, how deviations from expected
    behavior are detected, what automated safety systems are in place, and
    how the participant will be protected in the event of a system malfunction.

    \item A clear statement that the participant has the right to request
    human-only care at any point during the trial, including during active
    robotic procedures, and that exercising this right will not affect the
    participant's eligibility for continued trial participation or standard
    medical care.

    \item A description of the emergency stop procedures available to the
    participant, including how to verbally signal distress, how to activate
    patient-accessible stop mechanisms where available, and the expected
    response time from clinical staff.

    \item A statement identifying the Principal Investigator, Physical AI
    Safety Officer, and patient advocate who can be contacted with questions
    or concerns about Physical AI interactions at any time.
\end{itemize}

\subsubsection{Data Collection and Privacy Elements}

All Physical AI trial consent forms must include the following data-related
elements:

\begin{itemize}
    \item A comprehensive list of all data types collected by Physical AI
    systems, including but not limited to: robot telemetry data (position,
    force, speed, torque, temperature), video recordings of patient-robot
    interactions, audio recordings if applicable, sensor data from
    patient-mounted or room-mounted sensors, AI decision logs documenting
    automated recommendations, and biometric data collected through
    robot-integrated sensors.

    \item An explanation of how each data type is stored, who has access to
    it, how long it is retained, and under what circumstances it may be
    shared with parties outside the immediate clinical team.

    \item A statement that all Physical AI-generated data is protected under
    HIPAA and applicable state privacy laws (including SB 892 provisions for
    California sites), with specific protections for video and audio
    recordings that may capture identifiable patient images or voices.

    \item A description of the federated learning framework, if applicable,
    explaining that de-identified data may be used to improve AI models
    across trial sites without the participant's raw data leaving the local
    site.

    \item A statement of the participant's right to access all Physical AI
    data generated during their care, to receive copies of such data in a
    readable format, and to request correction of data errors.

    \item A description of cybersecurity protections for Physical AI data,
    including encryption standards, access controls, and incident response
    procedures in the event of a data breach.
\end{itemize}

\subsubsection{Tier-Specific Consent Elements}

In addition to the general and data elements, the following tier-specific
consent elements must be included based on the Physical AI systems involved
in the participant's care:

\begin{itemize}
    \item \textbf{Tier 2 additional elements:} A description of the specific
    physical interactions the participant will experience with each Tier 2
    robot, including what body parts will be contacted, the nature of the
    contact (positioning, support, imaging), the expected duration of contact,
    and any sensations the participant may experience (pressure, warmth,
    vibration). A description of the operator qualifications and the
    operator's role during the interaction. An explanation of the real-time
    monitoring systems active during contact and the thresholds that trigger
    automatic safety responses.

    \item \textbf{Tier 3 additional elements:} A detailed description of
    each surgical or invasive procedure involving Physical AI systems,
    including the specific role of the robot (instrument guidance, tissue
    manipulation, needle placement), the degree of autonomous operation
    permitted, and the surgeon's real-time control capabilities. A
    description of the emergency conversion plan detailing how the procedure
    transitions from robotic to manual if a system failure occurs. A
    statement of the risks specific to robotic surgical procedures, including
    mechanical failure during tissue interaction, sensor error leading to
    positional inaccuracy, and extended procedure time due to system
    recalibration. A description of the post-procedure Physical AI
    assessment, including the 72-hour follow-up schedule and the specific
    evaluations performed to detect robot-related complications.
\end{itemize}

\subsubsection{Ongoing Consent and Notification Elements}

The consent form must also address the ongoing consent requirements unique to
Physical AI trials:

\begin{itemize}
    \item A statement that Physical AI systems may be updated, replaced, or
    supplemented during the trial period, and that the participant will be
    notified in writing before any such change takes effect.

    \item A description of the specific changes that trigger re-consent,
    including introduction of new robot types, software updates that alter
    robot behavior or decision-making, changes to the degree of autonomy,
    and changes to safety monitoring procedures.

    \item A statement that the participant will receive periodic summary
    reports of Physical AI system performance at the trial site, including
    any adverse events that may affect the participant's risk assessment.

    \item An explanation that the participant may withdraw consent for
    Physical AI interactions at any time without withdrawing from the
    clinical trial itself, and a description of how the trial protocol
    accommodates continued participation with human-only care delivery.

    \item A statement that the participant's consent preferences will be
    documented in the electronic health record and communicated to all
    Physical AI systems through the MCP server infrastructure, ensuring
    that robots are automatically informed of any restrictions on their
    interaction with the participant.
\end{itemize}

\subsubsection{Comprehension Verification}

The adapted Part 50 requires that informed consent for Physical AI trials
include a comprehension verification step beyond the standard consent
signature. This verification must confirm that the participant understands:
the types of robots involved in their care and their general functions, the
participant's right to refuse any robotic interaction, how to activate
emergency stop procedures, what data is collected by Physical AI systems and
how it is protected, and whom to contact with questions or concerns about
Physical AI interactions. Comprehension may be verified through a brief
questionnaire, a teach-back conversation with the consent coordinator, or
another method approved by the IRB. The verification results must be
documented in the participant's trial record.
